% AUTOMATISCH ERZEUGT aus dem SysML-Modell mit tools/generate_docs.py.
% Nicht von Hand aendern: Inhalte im Modell pflegen, Layout im Kapitel.
SN-001 & STK-001 & Der Bediener moechte einen Transportauftrag mit moeglichst wenig Bedienaufwand ausloesen. & Must \\
SN-002 & STK-002 & Die Materialkiste soll am vorgesehenen Zielstandort ankommen. & Must \\
SN-003 & STK-003 & Der Betreiber moechte den Betriebszustand des Roboters jederzeit erkennen koennen. & Must \\
SN-004 & STK-004 & Personen im Fahrbereich duerfen durch den Roboter nicht gefaehrdet werden. & Must \\
SN-005 & STK-005 & Wartungspersonal benoetigt Diagnoseinformationen fuer Fehleranalyse und Instandhaltung. & Should \\
SN-006 & STK-006 & Der Roboter soll Hindernisse und geaenderte Umgebungen so behandeln, dass Betriebsablaeufe nicht unnoetig gestoert werden. & Must \\
SN-007 & STK-007 & Netzwerkschnittstellen, Updates und Datenhaltung sollen kontrollierbar und sicher sein. & Should \\
SN-008 & STK-009 & Architekturentscheidungen sollen nachvollziehbar dokumentiert und mit Anforderungen verknuepft sein. & Must \\
SN-009 & STK-011 & Der Roboter soll sich fuer die Demonstration technischer Zusammenhaenge im Unterricht eignen. & Must \\
SN-010 & STK-011 & Wichtige Systemzustaende und Ablaeufe sollen im Unterricht nachvollziehbar dargestellt werden koennen. & Should \\
